\section{Parallels in Aquinas's Other Works}\label{app:parallels}

Where Aquinas treats the same matter elsewhere: the \emph{Scriptum}
on the Sentences, the \emph{Contra Gentiles}, the disputed
questions \emph{De veritate}, \emph{De spiritualibus creaturis},
and \emph{De malo}, \emph{De ente et essentia}, and \emph{De
substantiis separatis}, in the order of the \emph{Summa}.

\begin{threestable}{Summa}{Parallel}{Agreement or difference}

I q.~22 a.~3; q.~103 a.~6 ad~1 & \emph{De sub. sep.} 3 & Difference of emphasis: \emph{De sub. sep.} 3 presents Platonic providence (higher providing for lower) as agreeing with Aristotle's; the \emph{Summa} rejects Plato's threefold division because God governs all things immediately in the design of government.\\
I q.~50 a.~1 & \emph{SCG} II.46, 49 & Agreement. \emph{SCG} II.46 argues the same need for intellectual creatures from assimilation to God (``ad perfectionem optimam universi, esse aliquas creaturas intellectuales''); II.49 argues that no intellectual substance is a body.\\
I q.~50 a.~2 & \emph{De spir. creat.} a.~1 & Same conclusion. The disputed question adds that if every potency is called matter, one may speak of matter and form in spirits ``ut non fiat vis in verbis'', but improperly.\\
I q.~50 a.~2 & \emph{De sub. sep.} cc. 5--8 & Fuller treatment of Avicebron. It diagnoses his two errors (taking logical composition as real; one meaning of potency, subject, recipient) and argues at length that spiritual and corporeal matter cannot be one (c. 7).\\
I q.~50 a.~2 & \emph{De spir. creat.} a.~1 co. & Agreement: the wide-sense matter/form name conceded but refused as improper.\\
I q.~50 a.~2 ad~3 & \emph{De ente} c. 4 (CT cap. 3); \emph{SCG} II.52--54 & Agreement. The composition of \emph{esse} and \emph{quod est}; ``potentia et actus … non tamen forma et materia nisi aequivoce'' (\emph{De ente}); \emph{SCG} II.54: composition from substance and \emph{esse} is not of the same kind as composition from matter and form.\\
I q.~50 a.~3 & \emph{SCG} II.92 & Agree: the separate substances exceed all material multitude; Aristotle's count is probable, not necessary; Maimonides' view excluded. \emph{SCG} adds the arguments from the order of the universe and from intelligible being.\\
I q.~50 a.~3 & \emph{De sub. sep.} 2 & Agree on number and on the probable character of Aristotle's argument; \emph{De sub. sep.} adds that Aristotle's way is surer but less sufficient than Plato's, because only the Platonists' demons explain possession and magic.\\
I q.~50 a.~4 & \emph{De spir. creat.} a.~8 & Same conclusion, with three reasons: the simple form cannot be multiplied; the order of the universe, which is \emph{per se} in the higher parts; and the perfection of the angelic nature. The \emph{Summa} gives chiefly the first.\\
I q.~50 a.~4 & \emph{De ente} c. 4 & Agreement: ``quot sunt ibi individua, tot sunt ibi species'', attributed to Avicenna.\\
I q.~50 a.~5 & \emph{SCG} II.55 & Agreement: the chapter title, ``Quod substantiae intellectuales sunt incorruptibiles''.\\
I q.~51 a.~1 (for the demons) & \emph{De malo} q.~16 a.~1 & Agreement: demons have no bodies naturally united; the question \latin{non multum refert ad fidei Christianae doctrinam}.\\
I q.~51 a.~1 ad~1 & \emph{De sub. sep.} 1 & Agree: the aerial bodies of demons belong to the Platonists' opinion, which Augustine uses without asserting.\\
I q.~51 a.~2 s.c. & \emph{De civ. Dei} XVI & Aquinas cites Augustine that angels appeared to Abraham in assumed bodies.\\
I q.~51 a.~3 ad~6 & \emph{De civ. Dei} XV.23 & Agreement: the sons of God are the sons of Seth; Aquinas quotes the chapter.\\
I q.~52 a.~1 & \emph{Quodl.} I q.~3 a.~1 & As Cajetan cites it (on I q.~52 a.~1): if ``operation'' means motion, the proposition is false; if any union of power, true.\\
I q.~54 a.~1 & \emph{SCG} II.97; \emph{De sub. sep.} c. 13 & Agreement: the separate substance's action is not its substance (c. 13 combats those who denied angels knowledge of singulars).\\
I q.~54 a.~4 & \emph{SCG} II.96; \emph{De ver.} q.~8 a.~8 & Agreement: no agent/possible intellect in separate substances, ``nisi forte aequivoce.''\\
I q.~54 a.~5 co. & \emph{De civ. Dei} XXI.10 & Aquinas: Augustine uses the opinion of bodies naturally united to angels and demons without asserting it; Augustine leaves it ``not … to be laboriously investigated''.\\
I q.~55 a.~2 & \emph{SCG} II.96; \emph{De ver.} q.~8 a.~9 & Agreement: knowledge not from sensible things; connatural species.\\
I q.~55 a.~2 co., ad~1 & \emph{De Gen. ad litt.} II.8.17 & Agreement: the knowledge of a thing's reason is first in the angelic mind.\\
I q.~55 a.~3 & \emph{De ver.} q.~8 a.~10 & Agreement: the higher angel knows through fewer, more universal forms.\\
I q.~56 aa.~1--3 & \emph{De ver.} q.~8 aa.~1--7 & Agreement in structure and determination.\\
I q.~56 a.~1 s.c.; a.~2 co. & \emph{De Gen. ad litt.} II.8.16 & Agreement: things proceed from the Word into the angelic mind and into their own natures.\\
I q.~57 a.~2 & \emph{De ver.} q.~8 a.~11 & Agreement; \emph{De veritate} names Avicenna as author of the universal-causes opinion the \emph{Summa} leaves unnamed.\\
I q.~57 a.~3 & \emph{De ver.} q.~8 a.~12; \emph{De malo} q.~16 a.~7 & Agreement: the future in itself is solius Dei; free-will futures conjectural even to angels and demons.\\
I q.~57 a.~4 & \emph{De ver.} q.~8 a.~13; \emph{De malo} q.~16 a.~8 & Agreement: thoughts not known directly, only per accidens through bodily signs and merit; \emph{De malo} quotes \emph{De div. daem.} and \emph{Retract.} II.30.\\
I q.~57 a.~4 co. & \emph{De div. daem.} 5.9; \emph{Retract.} II.30 & Agreement, with Augustine's own retraction quoted.\\
I q.~57 a.~5 obj.~1 & \emph{De Gen. ad litt.} V.19.38 & Used as objection (angels knew the mystery from the ages); Aquinas distinguishes.\\
I q.~58 aa.~1--5 & \emph{SCG} II.97, 101; \emph{De ver.} q.~8 aa.~14--15 & Agreement on non-discursion; difference of emphasis: \emph{SCG} II.97 says the separate intellect is ``semper intelligens actu'' outright, while \emph{ST} q.~58 a.~1 allows potency in the second mode for natural knowledge.\\
I q.~58 aa.~6--7 & \emph{De ver.} q.~8 aa.~16--17; \emph{De sub. sep.} cc. 19--20 & Agreement on the morning/evening division (\emph{De ver.} a.~16 gives the same ``ex parte rei cognitae / ex parte medii'' distinction); difference: \emph{De ver.} a.~17 restricts morning/evening, strictly taken, to the blessed angels' gratuitous knowledge and denies both to the demons, which the \emph{Summa} leaves implicit. \emph{De sub. sep.} cc. 19--20 marshals the same authorities (\emph{DN} 4; Gen. ad litt. VIII).\\
I q.~58 aa.~6--7 s.c. & \emph{De Gen. ad litt.} IV.22--24; \emph{De civ. Dei} XI.7, 29 & Agreement; Aquinas adds whether evening knowledge is by innate species (a.~7).\\
I q.~58 a.~1 s.c. & \emph{De Gen. ad litt.} II.8.17 & Agreement: the angels contemplate from their creation.\\
I q.~59 aa.~1--3 & \emph{SCG} II.46--48 & Agreement: will and free choice follow intellect in separate substances.\\
I q.~59 a.~4 ad~2 & \emph{De civ. Dei} IX.5 & Agreement: passions not said of angels as passions.\\
I q.~60 a.~5 obj.~4 & \emph{De civ. Dei} XII.9 & Objection from charity poured out in the angels.\\
I q.~61 a.~1 ad~1 & \emph{De civ. Dei} XI.9 (New Advent prints ``xi, 50'') & Agreement: angels designated by ``heaven'' or ``light''.\\
I q.~61 a.~2 & \emph{De potentia} q.~3 a.~17 & Agreement: the world (and angelic creation) is not eternal; faith, not demonstration.\\
I q.~61 a.~2 ad~2 & \emph{De Gen. ad litt.} VIII.20.39 & Agreement: God moves the spiritual creature through time.\\
I q.~61 a.~3 & \emph{De civ. Dei} XI.32--33 & Agreement: creation with the world more probable; contrary not erroneous; ``In the Son'' reading.\\
I q.~61 a.~3 & \emph{In II Sent.} d. 2 q.~1 a.~3 & Agreement: simultaneous creation more probable, other opinion not to be prejudged ``quia non est demonstratum, nec fide expressum''.\\
I q.~61 a.~3 & \emph{De potentia} q.~3 a.~18 & Agreement: ``Angeli simul cum creatura corporali sunt conditi; tamen sine alterius opinionis praeiudicio''; Jerome's report explained ``secundum opinionem antiquorum doctorum'' (ad~7).\\
I q.~61 a.~4 obj.~2, ad~2 & \emph{De Gen. ad litt.} III.10.14 & Aquinas reads the ``upper atmosphere'' as said of the angels who sinned.\\
I q.~62 a.~3 & \emph{In II Sent.} d. 4 q.~1 a.~3 & Development: the commentary calls the natural-state opinion ``communior'' while preferring created-in-grace; the \emph{Summa} drops the ``communior'' framing and judges created-in-grace ``probabilius et magis dictis sanctorum consonum''.\\
I q.~62 a.~3 obj.~1, s.c., ad~1 & \emph{De civ. Dei} XII.9; \emph{De Gen. ad litt.} II.8.16 & Agreement: created in grace (more probable); formation prior by nature, not time; Augustine leaves open unequal grace or more help to the good.\\
I q.~63 aa.~2--3 & \emph{In II Sent.} d. 5 q.~1 aa.~2--3 & Agreement in substance: pride as the first sin, arising from the sight of his own beauty; the Sentences reject equality simpliciter and place the desire in the mode of having beatitude (by his natural powers, without merit), as a.~3 does.\\
I q.~63 a.~1 & \emph{In II Sent.} d. 5 q.~1 a.~1 & Agreement: the angel sins by withdrawing consideration from one thing to another, without passion.\\
I q.~63 a.~1 ad~4 & \emph{De civ. Dei} XII.6--8 & Parallel: sin as a good sought without its rule.\\
I q.~63 a.~2 s.c., co. & \emph{De civ. Dei} XIV.3; \emph{De Gen. ad litt.} XI.14.18 & Agreement: pride and envy only, envy after pride.\\
I q.~63 a.~3 & \emph{De malo} q.~16 a.~3 & Agreement; \emph{De malo} adds that the sin lies in the supernatural order (\latin{sine Deo gratiam conferente}).\\
I q.~63 a.~3 & \emph{De malo} q.~16 a.~3 co. & Agreement, sharpened: beatitude ``per virtutem suae naturae... sine Deo gratiam conferente''.\\
I q.~63 a.~4 & \emph{De malo} q.~16 a.~2 & Agreement: evil by will, not nature.\\
I q.~63 a.~5 & \emph{In II Sent.} d. 3 q.~2 a.~1 & Agreement: the first-instant opinion ``vana et erronea et falsa'', condemned by the masters; the two-changes argument judged not cogent, as in the \emph{Summa}.\\
I q.~63 a.~5 & \emph{De malo} q.~16 a.~4 & Same conclusion, different reason: \emph{De malo} sets aside the argument from the nature as God made it and argues from the order of the angel's acts (natural self-knowledge in the first instant).\\
I q.~63 a.~5 & \emph{De malo} q.~16 a.~4 co. & Agreement, with history: the first-instant position ``reprobata fuit ab omnibus magistris tunc Parisiis legentibus''; and Aquinas's own reason (first instant = natural self-knowledge).\\
I q.~63 a.~5 co., ad~1, ad~2, ad~4 & \emph{De civ. Dei} XI.13, 15, 19 (New Advent prints ``xi, 15'' for XI.19) & Difference: the opinion of a devil wicked from the first instant, which Augustine does not count Manichaean and toward which \emph{De Gen. ad litt.} XI.16.21, 23.30 incline, is ``reasonably rejected by the masters as erroneous''.\\
I q.~63 a.~6 s.c., ad~4 & \emph{De civ. Dei} XI.15; \emph{De Gen. ad litt.} IV.24.41 & Aquinas determines the fall immediately after the first instant.\\
I q.~63 a.~7 & \emph{In II Sent.} d. 6 q.~1 a.~1 & Agreement: Gregory determines the highest of all, the common and probable opinion; the \emph{Summa} adds Damascene's alternative as not contrary to faith.\\
I q.~63 a.~7 co. & \emph{De Gen. ad litt.} XI.17.22, 19.25--26, 26.33; III.10.14 & The lower-order opinion (Damascene) against Gregory's; Augustine reports it without adopting it.\\
I q.~63 a.~8 & \emph{In II Sent.} d. 6 q.~1 a.~2 & Agreement: priority of causality, not of duration; the Sentences already cite Damascene's fall-as-death.\\
I q.~63 a.~9 ad~3 & \emph{Enchir.} 29; \emph{De civ. Dei} XXII.1 & Agreement: men taken into every order to supply the ruin.\\
I q.~64 a.~1 & \emph{In II Sent.} d. 7 q.~2 a.~1; \emph{De malo} q.~16 a.~6 & Agreement: natural intellectual light remains; \emph{De malo} adds that false opinion is possible only about what exceeds natural knowledge.\\
I q.~64 a.~1 obj.~3--5, co., ad~4--5 & \emph{De Gen. ad litt.} IV.22, V.19; \emph{De civ. Dei} IX.21, XI.33 & Agreement: Christ known to demons by temporal effects.\\
I q.~64 a.~2 & \emph{In II Sent.} d. 7 q.~1 a.~2 & Difference of argument: the Sentences reject the formula ``mobile before choice, immobile after'' as presupposing what is sought and ground obstinacy in the end of the way (fall as death); the \emph{Summa} adopts the formula and grounds it in the immobility of angelic apprehension.\\
I q.~64 a.~2 & \emph{De malo} q.~16 a.~5 & Synthesis: intrinsic immobility of the angelic will joined to the extrinsic end of the state of wayfarer; \latin{vis primae electionis}.\\
I q.~64 a.~4 & \emph{In II Sent.} d. 6 q.~1 a.~3 & Agreement: the air as penal and for our exercise, hell for their fault.\\
I q.~64 a.~4 s.c. & \emph{De Gen. ad litt.} III.10.15 & Agreement: the dark air a prison until judgment.\\
I q.~106 a.~1 & \emph{De ver.} q.~9 a.~1 & Agreement; \emph{De ver.} adds the analogy of bodily light (making visible and strengthening sight) and the higher spirit as act of the lower.\\
I q.~106 a.~1 & \emph{In II Sent.} d. 11 q.~2 a.~2 & Agreement: the more limpidly one sees the essence, the more effects one knows in it.\\
I q.~106 a.~2 ad~1 & \emph{De ver.} q.~9 a.~3 & Agreement, fuller: the three acts all belong to reception of knowledge; terminus a quo in angels only nescience.\\
I q.~106 a.~2 ad~2 & \emph{In II Sent.} d. 9 q.~1 a.~4 ad~5 & Agreement: Seraphim do not infuse charity but enlighten about what pertains to it.\\
I q.~106 a.~3 & \emph{De ver.} q.~9 a.~2 & Agreement; \emph{De ver.} records the contrary opinion and gives three reasons against it.\\
I q.~106 a.~4 ad~2 & \emph{In II Sent.} d. 11 q.~2 a.~4 & Agreement, finer: substance of the mystery known from the beginning, circumstances learned in fulfilment; Jerome, Augustine, Dionysius set side by side.\\
I q.~107 a.~1 & \emph{De ver.} q.~9 a.~4 & Agreement (three states of the intelligible form; will orders the concept).\\
I q.~107 a.~1 & \emph{In II Sent.} d. 11 q.~2 a.~3 & Difference: the Sentences commentary places speech in a naturally knowable sign coordinated with the concept (intellectual signs); \emph{ST} and \emph{De ver.} need no sign.\\
I q.~107 a.~2 & \emph{De ver.} q.~9 a.~5 & Agreement; \emph{De ver.} defines speech as making present what was absent, illumination as strengthening.\\
I q.~107 a.~3 & \emph{In II Sent.} d. 11 q.~2 a.~5 & Complement: the ``combat'' of Dan 10:13 as contrary merits referred to God; concord in receiving the divine illumination.\\
I q.~107 a.~4 & \emph{De ver.} q.~9 a.~6 & Agreement; \emph{De ver.} ad~2 restricts ``an angel works where he is'' to action on bodies.\\
I q.~107 a.~5 & \emph{De ver.} q.~9 a.~7 & Agreement.\\
I q.~108 aa.~5--6 & \emph{Enchir.} 58 & Difference: Augustine declines to distinguish the orders; Aquinas assigns them after Dionysius and Gregory.\\
I q.~108 a.~1 & \emph{In II Sent.} d. 9 q.~1 aa.~1, 3 & Agreement; etymology hieron + archon; royal-court image; Trinity reflected in three by three.\\
I q.~108 a.~1 & \emph{SCG} III.80 & Agreement: grades of knowing the order of providence (end, form, disposition; universal causes; particular causes).\\
I q.~108 a.~1 co. & \emph{De civ. Dei} XII.1 & Agreement: one society of good angels and men.\\
I q.~108 a.~3 & \emph{In II Sent.} d. 9 q.~1 a.~5 & Agreement: no two angels equal in the whole hierarchy; beauty in ordered grades.\\
I q.~108 a.~4 & \emph{In II Sent.} d. 9 q.~1 a.~7 & Agreement: grace formal, nature material and dispositive; orders in a manner from creation.\\
I q.~108 a.~4 obj.~2 & \emph{De civ. Dei} XII.9 & Charity poured into the angels (objection).\\
I q.~108 a.~5 ad~1 & \emph{Enarr. in Ps.} 103, s. 1.15; \emph{De civ. Dei} XV.23 & Parallel etymology: \latin{Angelus nuntius dicitur}.\\
I q.~108 a.~5 ad~1 & \emph{In II Sent.} d. 9 q.~1 a.~4 ad~2 & Agreement: ``angel'' name of nature as manifestative, name of office as exterior ministry (Gregory).\\
I q.~108 a.~6 & \emph{SCG} III.80 & Agreement: both orders have apostolic authority and differ little.\\
I q.~108 a.~6 & \emph{In II Sent.} d. 9 q.~1 a.~3 co., ad~7 & Difference: Dionysius's ordering ``more reasonable'', Gregory ``numbers rather than orders''; orders of each hierarchy reasoned by acts of fruition, order of peace, and limited power.\\
I q.~108 a.~6 co. & \emph{De Trin.} III.4.9 (New Advent prints ``\emph{De Trin.} ii''; Leonine ``III de Trin.'') & Agreement; used for Gregory's arrangement.\\
I q.~108 a.~7 & \emph{In II Sent.} d. 11 q.~2 a.~6 & Agreement: essentials of hierarchy remain; prelacy for leading wayfarers ceases.\\
I q.~108 a.~8 & \emph{In II Sent.} d. 9 q.~1 a.~8 & Agreement, fuller: three positions (Eustratius; virgins or perfect; all elect by merits), the Blessed Virgin above all; number known to God alone.\\
I q.~108 a.~8 co. & \emph{De civ. Dei} XII.9; XXII.29; \emph{Enchir.} 63 & Agreement: one society; men equal to angels by grace.\\
I q.~109 a.~1 obj.~3 & \emph{In II Sent.} d. 6 q.~1 a.~1 & Complement: Gregory and common opinion that the first sinner was highest; Augustine left it undetermined.\\
I q.~109 a.~2 & \emph{In II Sent.} d. 6 q.~1 a.~4 & Agreement, fuller: order from nature, divine wisdom, and their wickedness; faith among robbers; heavier torment for higher office.\\
I q.~109 a.~4 s.c. & \emph{De Trin.} III.4.9 (New Advent ``ii''; Leonine ``III'') & Agreement: the sinful spirit ruled by the just.\\
I q.~110 a.~1 & \emph{SCG} III.78 nn. 1--6 & Agreement; \emph{SCG} adds that intellectual creatures share both the ordering and the execution of providence and act by free will.\\
I q.~110 a.~1 ad~2--3 & \emph{De sub. sep.} 2--4 & Agree: Aristotle set movers over the heavens only and never mentions demons; the Platonists set spiritual substances over bodies, as the holy doctors do.\\
I q.~110 a.~1 ad~3 & \emph{SCG} III.80 n. 11 & Agreement: the Virtues move the heavens and execute works beyond nature, with Gregory's sentence.\\
I q.~110 a.~1 s.c., co. & \emph{De Trin.} III.4.9; \emph{De div. qq. 83} q.~79.1 & Agreement.\\
I q.~110 a.~2 & \emph{De pot.} q.~6 a.~3; \emph{De malo} q.~16 a.~9 & Agreement; \emph{De malo} grounds it in the order by which the lowest are moved by the highest through intermediaries.\\
I q.~110 a.~2 s.c. & \emph{De Trin.} III.8.13 & Agreement: matter obeys God, not the rebel angels.\\
I q.~110 a.~3 & \emph{De malo} q.~16 a.~10; \emph{De pot.} q.~6 a.~3 & Agreement; \emph{De malo} adds that local motion is the least change and can come from a remote agent; \emph{De pot.} calls such works per modum artis. \emph{De malo} a.~10 leaves open whether the demons were of the lower angels over the earthly order or fell from the highest (Gregory).\\
I q.~110 a.~3 s.c. & \emph{De Trin.} III.8--9 & Agreement: seeds applied by angels.\\
I q.~110 a.~4 & \emph{SCG} III.103; \emph{De pot.} q.~6 aa.~3--4 & Agreement; \emph{SCG} answers Avicenna at length and allows miracles by spiritual substances acting in divine power; \emph{De pot.} names three modes of angelic share in miracles.\\
I q.~110 a.~4 obj.~2 & \emph{De div. qq. 83} q.~79 & Objection.\\
I q.~111 aa.~3--4 & \emph{De malo} q.~16 a.~11 & Agreement; \emph{De malo} tests Augustine's three accounts (\emph{De Gen. ad litt.} XII) and admits only the local change of what pre-exists in the body.\\
I q.~111 a.~1 & \emph{De malo} q.~16 a.~12 & Agreement; \emph{De malo} adds that the good angels strengthen the human light and the demons do not, and that demons arrange images toward pride or insoluble doubt.\\
I q.~111 a.~2 ad~2; q.~114 a.~2 ad~2 & \emph{De malo} q.~16 a.~8; \emph{ST} I q.~57 a.~4 & Agreement; \emph{De malo} explains why the act of the will is hidden to all but God and the one willing, while species in the intellect may be known by a higher intellect.\\
I q.~112 a.~4 & \emph{SCG} III.80 nn. 5--19 & Agreement on the offices of the orders; \emph{SCG} follows Dionysius's order and reports Gregory's.\\
I q.~112 a.~4 ad~2 & \emph{In De div. nom.} prooem. & Both accept a Platonic principle only where it is consonant with faith (fewer orders nearer the first principle; the doctrine of the first principle), rejecting separate forms.\\
I q.~113 a.~1 & \emph{In II Sent.} d. 11 q.~1 a.~1 & Agreement; the Sentences adds the reasons of fittingness (dignity of angels as ``duces hominum reductionis in Deum''), angels prepare for grace, carry prayers, do not compel, rarely appear.\\
I q.~113 a.~1 (with q.~57 a.~2) & \emph{De veritate} q.~8 a.~11 & Agreement: guardianship requires knowledge of singulars.\\
I q.~113 a.~2 & Lombard, \emph{Sent.} II d. 11 c. 1 (no Thomist article) & Difference with the Master: one angel for each man versus one angel for many.\\
I q.~113 a.~3 & \emph{In II Sent.} d. 11 q.~1 a.~2; \emph{SCG} III.80 & Agreement; Sentences s.c. cites Bernard, \emph{De consid.}; \emph{SCG} assigns objects of Principalities, Archangels, Angels.\\
I q.~113 a.~4 & \emph{In II Sent.} d. 11 q.~1 a.~3 & Agreement (Adam, Antichrist, Christ); Sentences adds ``lex communis propter unum mutari non debet'' and answers \emph{CH} 4.4 on Christ.\\
I q.~113 a.~5 & \emph{In II Sent.} d. 11 q.~1 a.~3 co. and ad~3 & Difference: Sentences gives a guardian from the infusion of the rational soul; \emph{Summa} gives the mother's angel in the womb and a guardian at birth.\\
I q.~113 a.~6 & \emph{In II Sent.} d. 11 q.~1 a.~4 & Agreement; Sentences adds the angels departing at the hour of death, the lasting impression of one devout prayer, and no one on the way despaired of.\\
I q.~113 a.~7 & \emph{In II Sent.} d. 11 q.~1 a.~5 & Agreement; Sentences adds impassibility, joy in their own good works, antecedent and consequent will.\\
I q.~113 a.~8 & \emph{In II Sent.} d. 11 q.~2 a.~5 & Agreement with Gregory; Sentences reports Jerome as holding an evil prince and adds the good the Jews did in Persia.\\
I q.~114 a.~1 & \emph{De pot.} q.~6 a.~10 & Related: the demons are coerced by God and, in God's power, by saints and angels; Daemones arcere belongs to the order of Powers.\\
I q.~114 a.~4 & \emph{De pot.} q.~6 a.~5 & Agreement; \emph{De pot.} gives the reason: a miracle is divine testimony, and God will not witness to the demons' falsehood.\\
I q.~114 a.~4 ad~2, ad~3 & \emph{De Trin.} III.8--9; \emph{De div. qq. 83} q.~79.4 & Agreement.\\
\end{threestable}
